\documentclass[10pt,a4paper]{amsart}
\usepackage[margin=19mm]{geometry}
\usepackage{amsmath,amssymb,mathtools}
\usepackage[scr=rsfs,cal=euler]{mathalfa}
\usepackage{xfrac}
\usepackage[unicode,hidelinks,bookmarks=false]{hyperref}
\newcommand{\ie}{i.e.\ }
\newcommand{\cf}{cf.\ }
\newcommand{\resp}{respectively\ }
\newcommand{\Subsection}{\subsection}
\providecommand{\nobreakdash}{\nobreak\hbox{-}}
\setlength{\emergencystretch}{2em}
\multlinegap0pt
\usepackage{amssymb}
\usepackage[shortlabels]{enumitem}
\newtheorem{proposition}{Proposition}
\title{A Linear-Size Block-Partition Fibonacci Encoding\\for G\"odel Numbering}
\author{Zolt\'an S\'ostai}
\date{Source: arXiv:2603.25307v1 (CC BY 4.0)}
\begin{document}
\maketitle

\noindent\emph{Source and licence.} A Linear-Size Block-Partition Fibonacci Encoding\\for G\"odel Numbering. Version arXiv:2603.25307v1 (CC BY 4.0), distributed by arXiv under the Creative Commons Attribution 4.0 International licence, http://creativecommons.org/licenses/by/4.0/, as recorded in the arXiv licence metadata for this exact version and shown on the abstract page https://arxiv.org/abs/2603.25307v1. Full licence notice: \texttt{licenses/arxiv-2603.25307-CC-BY-4.0.txt}.\par\smallskip

\noindent\textit{Editorial scope.} Counted unit: the article's proposition prop:blowup (source0.tex lines 326--328). The article's complete original proof of it is delivered with it (source0.tex lines 330--343, one \texttt{proof} environment of 14 lines). No mathematics is written, repaired or re-derived: the statement and the proof are the article's own lines, in the article's own notation, and no retained line is substituted or re-flowed.  No further source-local context is needed: every cross-reference of every retained line resolves inside the two delivered spans.  Nothing was omitted from any retained span, so the package carries no omission record.  No retained line contains a citation, so the wrapper carries no bibliography and no bibliography entry.  No retained line contains a figure environment, an image-inclusion command or drawing code, so no image is delivered and the package declares no asset dependency.  Editorial formatting only, in addition: an AMS \texttt{amsart} preamble that restates the article's own declarations and macros used by the retained lines, namely the environments \texttt{proposition} on separate counters, together with the standard packages those lines need.  The retained environments print in this document as Proposition 1 in order of appearance, whereas the article prints the same items with the numbering of its own full document.  The complete native source of the article as distributed in the arXiv e-print for this exact version is bundled unchanged as \texttt{sources/197198/source0.tex}.
\begin{proposition}[Exponential growth]\label{prop:blowup}
Let codes be drawn from $\{1,\ldots,k\}$. Right-nested encoding of a length-$m$ sequence via Rosko's carryless pairing produces codes whose worst-case digit count is $\Theta(2^m)$.
\end{proposition}

\begin{proof}
In Rosko's pairing $\pi(x,y)$, the output's Zeckendorf support occupies indices up to $B(x) + 2\,R(y) - 1$, where $B(x) = 2\,r(x)$ is the delimiter, $r(x) = \min\{e : F_e > x\}$ is the rank, and $R(y) = \max Z(y)$. Since all symbol codes are positive, $Z(a_i)$ is nonempty and $R(a_i)$ is well-defined for every symbol.

Over the fixed symbol set $\{1,\ldots,k\}$, define
\[
b_{\min} = 2 \cdot \min\bigl\{\, r(a) : a \in \{1,\ldots,k\} \,\bigr\}, \qquad
b_{\max} = 2 \cdot \max\bigl\{\, r(a) : a \in \{1,\ldots,k\} \,\bigr\}.
\]
At each nesting step, the odd band determines the growth:
\[
2\,M_j + b_{\min} - 1 \;\le\; M_{j+1} \;\le\; 2\,M_j + b_{\max} - 1.
\]
Both bounds are linear recurrences of the form $M_{j+1} = 2M_j + C$, solving to $A \cdot 2^j - C$ for a constant~$A$ determined by initial conditions. Therefore $M_j = \Theta(2^j)$. Since the digit count is $\Theta(M_j \cdot \log_{10}\varphi)$, the encoded number has $\Theta(2^j)$ digits. A length-$m$ sequence requires $m-1$ pairings, giving worst-case $\Theta(2^m)$ digit growth.
\end{proof}

\end{document}
